\subsubsection{General Free-flow Boundary Condition}
\label{SEC:BC:Farfield:General}

This boundary condition serves as a generalized inflow\slash outflow for both
subsonic and supersonic flow regimes, essentially combining the general inflow
(\ref{SEC:BC:Inflow:General}) and general outflow (\ref{SEC:BC:Outflow:General})
boundary conditions into a single option.
%
This boundary condition takes the velocity from the left\slash interior state,
$\vec{u}_\subL$, and evaluates the magnitude of that velocity normal to the
boundary face, $v_{n,\subL}$, as
%
\[
  v_{n,\subL} = \vec{u}_\subL\cdot\normal
\]
%
where \normal{} is the unit normal for the boundary face that points outside the
computational domain.
%
The interior static temperature $T_\subL$ and $v_{n,\subL}$ can then be used to
compute an interior Mach number normal to the inflow boundary, $M_{n,\subL}$, as
%
\[
  M_{n,\subL} = \frac{v_{n,\subL}}{\sqrt{\gamma \rgas T_\subL}}
\]
%
From here, there are four options based on the magnitude and sign of
$M_{n,\subL}$:
%
\begin{itemize}
  %
  \item $M_{n,\subL}\geq +1$\quad (\textit{Supersonic Outflow})

    The interior flow is supersonic and exiting the domain.
    %
    The right state is evaluated the same way as the supersonic outflow boundary
    condition (\ref{SEC:BC:Outflow:Supersonic}) by just fixing it to the
    interior state, \thatis \[ \cvVecR = \cvVecL \]
  %
  \item $M_{n,\subL}\leq -1$\quad (\textit{Supersonic Inflow})

    The interior flow is supersonic because the magnitude of the Mach number is
    greater than or equal to 1. It is entering the domain since a negative value
    indicates the flow is aligned in the opposite direction as the outward
    facing normal.
    %
    The right state is evaluated the same way as the supersonic inflow boundary
    condition (\ref{SEC:BC:Inflow:Supersonic}) by just fixing it to the boundary
    condition state, \thatis \[ \cvVecR = \cvVecBC \]
  %
  \item $0 \leq M_{n,\subL} < 1$\quad (\textit{Subsonic Outflow})

    The interior flow is subsonic and exiting the domain.
    %
    The right state is evaluated based on the subsonic outflow boundary
    condition (\ref{SEC:BC:Outflow:Subsonic}), except there are no safe guards
    in place for supersonic or reversed flows since the other three conditional
    statements of this bulleted list explicitly deal with those situations for
    this boundary condition.
  %
  \item $-1 < M_{n,\subL} < 0$\quad (\textit{Subsonic Inflow})

    The interior flow is subsonic because the magnitude of the Mach number is
    less than 1; it is entering the domain since a negative value indicates the
    flow is aligned in the opposite direction as the outward facing normal.
    %
    The right state is evaluated based on the subsonic inflow boundary condition
    with $\inptext{Sub\_Inflow\_Method}=4$ (\ref{SEC:BC:Inflow:Subsonic:4}).
    %
    It should be noted that this boundary condition will always operate as if
    $\inptext{Sub\_Inflow\_Method}=4$ when the interior Mach number indicates a
    subsonic inflow condition; the value of the input
    \inptext{Sub\_Inflow\_Method} has no control over the subsonic inflow
    method used for this boundary condition and is therefore ignored.
  %
\end{itemize}

%%%%%%%%%%%%%%%%%%%%%%%%%%%%%%%%%%%%%%%%%%%%%%%%%%%%%%%%%%%%%%%%%%%%%%%%%%%%%%%%

\devNotes{Need to fix subsonic inflow angle-of-attack bug.}

%###############################################################################

